\documentclass[11pt]{article}
\usepackage[a4paper,margin=25mm]{geometry}
\usepackage{amsmath,amssymb,amsthm}
\usepackage[british]{babel}
\usepackage[round,authoryear]{natbib}
\usepackage{url}
\usepackage[hidelinks]{hyperref}
\newtheorem{theorem}{Theorem}
\newtheorem{proposition}[theorem]{Proposition}
\newtheorem{lemma}[theorem]{Lemma}
\newcommand{\E}{\mathbb E}
\newcommand{\Prb}{\mathbb P}
\title{Supporting derivation for a shared-threshold positive pricing correction}
\author{Amanjeet Singh}
\date{10 October 2026}
\begin{document}
\maketitle
\noindent Keywords: quantum derivative pricing; positive threshold measures;
Gaussian smoothing; reversible interval predicates; financial error certificates.

\section{Target and scope}

I retain the established classical risk-neutral model of the main paper.
The independent common Gaussian factor has standard deviation $\lambda>0$;
$Y$ denotes the remaining factors. At a given $y$, the geometric log-level
is $m=m_G(y)$ and the arithmetic log-level is
$m_A(y)=m+\Delta(y)$, where $\Delta=\log T\geq0$ by the model's positive
weights and exact weighted-centred drift and loading conditions.
Geometric conditioning and conditional payoff smoothing have published
financial antecedents \citep{curran,bayer}; the retained source ledger
records the versions previously consulted.
The target correction is $R_1(y)=P_\lambda(m+\Delta)-P_\lambda(m)$.
All financial coefficients, the outer factor law and the no-arbitrage
interpretation are unchanged. I alter only the coherent representation of
that correction. Positive layer-cake representations, common random
numbers, Gaussian smoothing and reversible comparisons are established
mechanisms. The new calculation below is a finite implementation and
complete financial certificate for their composition, not a claim of
generic theoretical priority.

\section{A positive finite-threshold identity}

\begin{theorem}[Shared-threshold correction]
Let $\nu$ be a positive finite measure on $\mathbb R$ of mass
$0<H<\infty$, and define
\[
 P_{\lambda,\nu}(m)=\int_{\mathbb R}
      \Phi((m-x)/\lambda)\,\nu(dx).
\]
Let $X$ have law $\nu/H$ and let $N$ be an independent standard normal
variable. For every finite $m\in\mathbb R$ and $\Delta\geq0$,
\[
 P_{\lambda,\nu}(m+\Delta)-P_{\lambda,\nu}(m)
   =H\Prb\{m<X-\lambda N\leq m+\Delta\}.
 \tag{S1}\label{eq:shared-identity}
\]
The represented score is in $[0,H]$ for every draw. This includes positive
finite mixtures of cash digitals and bounded increasing payoffs admitting
the stated positive threshold measure. Negative measures and an
infinite-mass representation of an uncapped vanilla payoff are excluded.
\end{theorem}
\begin{proof}
Conditional on $X=x$, Gaussian symmetry gives
\[
 \Prb\{x-\lambda N\leq m\}
 =\Prb\{N\geq(x-m)/\lambda\}
 =\Phi((m-x)/\lambda).
\]
Integration against the probability law $\nu/H$ gives the distribution
function of $X-\lambda N$. Subtract its values at $m+\Delta$ and $m$.
Because $\Delta\geq0$, the two events are nested. Their difference is
precisely the interval indicator in \eqref{eq:shared-identity}, multiplied
by $H$. The independent Gaussian makes its distribution continuous even
when $\nu$ has atoms, so continuous endpoint conventions are equivalent.
No independence between the two endpoint events is introduced.
\end{proof}

The measure theorem is an expectation identity, not a free state-access
assumption. A finite implementation for a different $\nu$ must construct
and charge a certified approximation to $\nu/H$, including any label
weights, reversible value map, preparation inverse and probability bias.
Only the stated uniform-strike construction is implemented in this study.

For the cap $F(s)=\min\{(s-K)^+,H\}$, the elementary threshold identity
$F(s)=\int_0^H\mathbf 1\{s>K+u\}\,du$ holds separately on
$s\leq K$, $K<s<K+H$ and $s\geq K+H$.
Substituting the conditional lognormal stock level $Ke^{m+\lambda N}$
and integrating the nonnegative bounded indicator gives
\[
 P_\lambda(m)=H\int_0^1
  \Phi\!\left(\frac{m-a(s)}{\lambda}\right)ds,
 \qquad a(s)=\log(1+Hs/K).
\]
Thus $X=a(S)$ for an independent uniform $S\in[0,1]$ implements the
cap case of \eqref{eq:shared-identity}. One shared threshold draw gives
a nonnegative correction with normaliser $H$, without evaluating two
conditional-price tables or subtracting two approximated prices.

The alternative elementary identity
$R_1=\Delta\int_0^1P_\lambda'(m+t\Delta)dt$ is valid, but a sampled
integrand is bounded only by $\Delta_*\|P_\lambda'\|_\infty$ under
that generic bound. Clipping it at $H$ can change its expectation. The
interval representation obtains $H$ directly from a Boolean event.

\section{Inner discretisation and exact confidence-independent errors}

Write $\phi$ for the standard normal density and choose certified
$\nu_0\geq\|\phi\|_\infty$ and
$\nu_1\geq\|\phi'\|_\infty$. These constants are separate from the
threshold measure $\nu$.

\begin{lemma}[Midpoint strike law]
Let $M_S$ uniform strike cells have midpoints $s_j=(j+1/2)/M_S$.
For all $m\in\mathbb R$ and $\Delta\geq0$, replacing the uniform
$S$ in the ideal correction by their equal-weight midpoint law has
absolute monetary error at most
\[
 \beta_S=\frac{H(H/K)^2}{12M_S^2}
       \left(\frac{\nu_1}{\lambda^2}
                         +\frac{\nu_0}{\lambda}\right).
 \tag{S2}\label{eq:strike-bound}
\]
\end{lemma}
\begin{proof}
Put $c=H/K$ and
$h(s)=\Phi((m+\Delta-a(s))/\lambda)-\Phi((m-a(s))/\lambda)$.
For either endpoint $v\in\{m,m+\Delta\}$, differentiating twice gives
\[
 \frac{d^2}{ds^2}\Phi((v-a(s))/\lambda)
 =\frac{a'(s)^2}{\lambda^2}\phi'((v-a(s))/\lambda)
   -\frac{a''(s)}{\lambda}\phi((v-a(s))/\lambda).
\]
Here $a'=c/(1+cs)$ and $a''=-c^2/(1+cs)^2$. The triangle inequality
therefore gives
$\|h''\|_\infty\leq2c^2(\nu_1/\lambda^2+\nu_0/\lambda)$.
On a cell of width $w=1/M_S$ with midpoint $s_j$, Taylor's theorem
and $|h''|\leq C$ imply
$|h(s)-h(s_j)-h'(s_j)(s-s_j)|\leq C(s-s_j)^2/2$.
The linear term integrates to zero and
$\int_{-w/2}^{w/2}t^2dt=w^3/12$. Summing all $M_S$ cells gives
$C/(24M_S^2)$ for the probability error. Multiplication by $H$ proves
\eqref{eq:strike-bound}.
\end{proof}

\begin{lemma}[Rounded truncated-normal midpoint law]
Let $L>0$, $a_N=\Phi(L)-\Phi(-L)$, and let $n_j$ be the $M_N$
equal-CDF-cell midpoints of the standard normal law conditional on
$[-L,L]$. Suppose actual dyadic nodes $\widehat n_j$ satisfy
$|\widehat n_j-n_j|\leq\delta_N$. Their equal-label empirical
distribution function $\widehat F_N$, with label multiplicities retained,
satisfies
\[
 \sup_x|\widehat F_N(x)-\Phi(x)|
 \leq D_N:=\frac{1-a_N}{2}+\frac{a_N}{2M_N}+\nu_0\delta_N.
 \tag{S3}\label{eq:normal-cdf}
\]
The same bound holds for $\widehat F_N(x-)$.
\end{lemma}
\begin{proof}
For the ideal midpoint law, let $F_T$ be the conditional normal CDF and
$F_M$ the midpoint empirical CDF. Equal-CDF midpoint placement gives
$|F_M(x)-F_T(x)|\leq1/(2M_N)$ on $[-L,L]$; at nodes the two
one-sided values attain the endpoint deviations. On this interval,
$\Phi(x)=(1-a_N)/2+a_NF_T(x)$. Consequently
\[
 F_M(x)-\Phi(x)
  =a_N(F_M(x)-F_T(x))+(1-a_N)(F_M(x)-1/2).
\]
Since $F_M\in[0,1]$, the claimed bound without the rounding term follows.
Outside $[-L,L]$, the deviation is at most $(1-a_N)/2$, so the same
bound holds globally. Coupling each rounded node with its labelled
ideal node gives
$F_M(x-\delta_N)\leq\widehat F_N(x)\leq F_M(x+\delta_N)$.
The mean value theorem and $\|\phi\|_\infty\leq\nu_0$ then add
at most $\nu_0\delta_N$ to the discrepancy from $\Phi(x)$.
Taking increasing sequences converging to $x$, continuity of $\Phi$
gives the identical bound for $\widehat F_N(x-)$. Duplicate rounded
nodes do not invalidate the coupling or empirical-CDF argument.
\end{proof}

\begin{proposition}[Arithmetic event sandwich]
At every exact actual outer financial node and ideal strike midpoint,
write $\ell=m+J$, $u=m+\Delta+J$ and
$z=a(s)+J-\lambda\widehat N$, where $\widehat N$ has the preceding
rounded midpoint law. Suppose the reversible body produces represented
values satisfying
\[
 |\widehat\ell-\ell|\leq e_\ell,\qquad
 |\widehat u-u|\leq e_u,\qquad
 |\widehat z-z|\leq e_z,
\]
and integer range guards exclude modular overflow. The implemented
Boolean event $\widehat\ell<\widehat z\leq\widehat u$ differs in
conditional probability from the ideal continuous-normal event by at most
\[
 2D_N+\frac{\nu_0}{\lambda}(e_\ell+e_u+2e_z).
 \tag{S4}\label{eq:sandwich}
\]
The quantile rounding already charged in $D_N$ is excluded from $e_z$;
the error of scaling the actual rounded node is included in $e_z$.
\end{proposition}
\begin{proof}
The implemented event is contained in
$\ell-e_\ell-e_z<z\leq u+e_u+e_z$ and contains
$\ell+e_\ell+e_z<z\leq u-e_u-e_z$ when the latter is nonempty.
An interval under $z=a(s)+J-\lambda\widehat N$ differs from its
continuous Gaussian counterpart by at most $2D_N$, because its probability
is a difference of two CDF or left-limit values and
\eqref{eq:normal-cdf} applies to each. Continuous $z$ has density bounded
by $\nu_0/\lambda$, so expanding or shrinking the endpoints adds or
removes probability at most
$\nu_0(e_\ell+e_u+2e_z)/\lambda$.
If the shrunken interval is empty, the ideal interval has width at most
$e_\ell+e_u+2e_z$; its entire continuous probability is then at most
the same density bound. The trivial lower bound zero gives the required
lower inequality in that case as well. Combining the two inequalities
proves \eqref{eq:sandwich}. No independence of endpoint arithmetic errors
is assumed. Finite atoms at strict endpoints are handled by the left-limit
CDF bound, rather than disregarded as if the discrete law were continuous.
\end{proof}

Thus a pointwise conditional probability allowance is
\[
 \eta=\min\left\{1,
   \beta_S/H+2D_N+
   \frac{\nu_0}{\lambda}(e_\ell+e_u+2e_z)\right\}.
 \tag{S5}\label{eq:eta}
\]
This bounds conditional expected flags against the conditional correction.
It does not equate a single Boolean outcome with a numerical payoff.
For the finite study $H=1$, $L=6$, $M_N=1024$ and $M_S=16$.
All outer scalar addresses remain at twelve bits; the small-range strike
logarithm uses six address bits. Endpoint and scaling errors are propagated
through actual scalar-table and unsigned arithmetic contracts, with the
existing outward certificates and overflow guards. The body certificate
treats normal coordinates as their actual rounded values; its zero
quantile error has this explicit meaning. The separate term
$\nu_0\delta_N$ in \eqref{eq:normal-cdf} restores the approximation to
the continuous inner normal. Multiplying again by an inner normal-box
mass would change the algorithm and invalidate this decoder.

\section{Coherent predicate and complete financial reconstruction}

The numeric body computes the common endpoint offset, nonnegative
represented log-ratio and shared threshold. I form the represented upper
endpoint by adding that same nonnegative log-ratio to the represented
lower endpoint. The range guards exclude overflow, so represented upper
is at least represented lower on every basis input. Two clean unsigned
comparisons produce
$b_\ell=\mathbf1\{\widehat\ell<\widehat z\}$ and
$b_u=\mathbf1\{\widehat u<\widehat z\}$.
The exact identity
\[
 b_\ell(1-b_u)
  =\mathbf1\{\widehat\ell<\widehat z\leq\widehat u\}
\]
holds also at ties and when the endpoints coincide. A negative control on
$b_u$ and a Toffoli XOR this predicate into an independently allocated
objective flag. I reverse the comparison work and then the numeric body.
The retained flag is excluded from those reversals. The result is the
exact permutation
$|j\rangle|f\rangle|0\rangle\mapsto
 |j\rangle|f\mathbin\oplus b(j)\rangle|0\rangle$ for either initial
flag value. Linearity gives the same coherent action on arbitrary
superpositions, including entanglement with an untouched reference.
The emitted exact Toffoli expansion supplies the required phases.

Hadamards prepare all outer financial, inner normal and strike labels.
Only those labels and the objective flag belong to the zero reflection;
all numeric and predicate workspace is clean in the invariant subspace.
There is no additional $q$-bit uniform probability-comparison register.
Every preparation and full adjoint, good/zero reflection, Fourier circuit,
readout and reset is charged in the retained complete QPE plans.
No free Gaussian state preparation, stochastic draw or conditional-price
oracle is assumed.
The canonical amplification and estimation mechanism is the established
one of \citet{brassard}; quantum Monte Carlo pricing has the published
financial antecedent of \citet{rebentrost}. I reuse their mechanism rather
than claim that a new financial flag creates a new generic estimator.

Let $c_0$ be the classical continuous coefficient, let $e_c$ bound its
chosen midpoint error, and let $\beta_Y$ bound the difference between the
original continuous correction and its mass-weighted actual outer-grid
correction. Let $[F_-,F_+]$ enclose the original outer Gaussian box mass,
with midpoint $F_0$ and radius $e_F$. If $p$ is the implemented mean
flag and the estimator returns $\widehat p\in[0,1]$, with probability
at least $1-\alpha$ its certified probability error is at most
$\varepsilon_p$. The price decoder is
$\widehat V=\widehat c_0+F_0H\widehat p$, and
\[
 |\widehat V-V|
 \leq e_c+\beta_Y+F_+H(\eta+\varepsilon_p)+e_FH.
 \tag{S6}\label{eq:complete-price}
\]
To prove this, insert the exact continuous coefficient, the mass-weighted
ideal outer-grid correction and its implemented conditional flag mean in
succession. Their discrepancies are bounded by $e_c$, $\beta_Y$ and
$F_+H\eta$. Replace $p$ by the certified estimator value and charge
$F_+H\varepsilon_p$. Finally, $H|\widehat p|\leq H$ bounds the error
from replacing the outer mass by its midpoint, giving $e_FH$.
Discounting multiplies the entire bound by $e^{-r_f\tau}$.
The existing complete outer-transfer certificate is retained because
the exact target $R_1$ has not changed.

An existing price-independent finite mean interval $[r_-,r_+]$ for
$\E_NR_1$ yields the implemented-flag interval
$[\max(0,r_-/H-\eta),\min(1,r_+/H+\eta)]$.
This follows by averaging the uniform conditional probability allowance.
The prior's construction uses the same analytic coefficient, continuous
second ratio moment and curvature certificate as the direct opponent;
no classical price answer or auxiliary outcome supplies it.

\section{Interpretation and falsification criteria}

The auxiliary score has conditional mean $R_1(Y)$, so the established
conditional variance identity gives
$\operatorname{Var}(H b)=\operatorname{Var}(R_1(Y))+
\E\operatorname{Var}(H b\mid Y)$ for the continuous ideal construction.
Classical integration of the auxiliaries recovers the old conditional
kernel and cannot increase its variance. I therefore claim no classical
variance advantage from this representation. A constructive quantum
improvement requires a lower completely charged preparation cost that
survives its larger certified deterministic allowance, its statistical
phase width and every confidence and reflection charge.
The retained finite study has this scope; it does not establish an
optimal oracle, general theoretical priority, physical runtime advantage
or failure of stronger classical pricing algorithms. Any arithmetic-domain
failure, omitted workspace inverse, extra uncharged mass factor,
uncontrolled endpoint atom or cost increase is a rejection under the
registered policy, rather than a reason to change the policy after outcomes.
The old numeric-payoff preparation also produces a Bernoulli good flag
through its independent uniform probability-comparison register. In the
continuous ideal limit, both preparations therefore have the same global
good probability $\E R_1/H$ and the same amplitude-amplification angle.
The conditional classical variance penalty does not imply an additional
quantum query penalty merely from the new auxiliary labels. Conversely,
the new representation does not reduce the ideal good-flag measurement
variance. Any query difference arises from its separately charged finite
errors and priors. The observed benefit here is cheaper emitted
preparation, with unchanged certified QPE phase sizes in the four cases.
An additional frozen numerical-rival pool varies its two conditional-price
table address widths independently over $\{6,8,10,12\}$ while granting
both coherent architectures and retaining every other financial input and
precision rule. All 64 outcomes, including twelve deterministic rejects,
are retained. Ten bits for each table gives the lowest complete numerical
T cap in all four cases. The fixed threshold preparation still improves
these stronger caps by approximately 12.33, 5.26, 6.26 and 2.95 per cent.
This is a finite-pool logical comparison, not a proof of global precision
optimality or a physical speed-up. The historical twelve-bit price tables
give larger percentages and are not used as the strongest tested rival.

\begin{thebibliography}{99}
\bibitem[Bayer et al.(2018)]{bayer}
Bayer, C., Siebenmorgen, M. and Tempone, R. (2018).
Smoothing the payoff for efficient computation of basket option prices.
\emph{Quantitative Finance}, 18(3), 491--505.
\url{https://doi.org/10.1080/14697688.2017.1308003}.
\bibitem[Brassard et al.(2002)]{brassard}
Brassard, G., H\o yer, P., Mosca, M. and Tapp, A. (2002).
Quantum amplitude amplification and estimation.
\emph{Contemporary Mathematics}, 305, 53--74.
\url{https://doi.org/10.1090/conm/305/05215}.
\bibitem[Curran(1994)]{curran}
Curran, M. (1994). Valuing Asian and portfolio options by conditioning
on the geometric mean price.
\emph{Management Science}, 40(12), 1705--1711.
\url{https://doi.org/10.1287/mnsc.40.12.1705}.
\bibitem[Rebentrost et al.(2018)]{rebentrost}
Rebentrost, P., Gupt, B. and Bromley, T. R. (2018).
Quantum computational finance: Monte Carlo pricing of financial derivatives.
\emph{Physical Review A}, 98, 022321.
\url{https://doi.org/10.1103/PhysRevA.98.022321}.
\end{thebibliography}
\end{document}
